\UnnumberedChapter{LỜI MỞ ĐẦU}

Nhận dạng biển báo giao thông là một thành phần quan trọng của các hệ thống hỗ trợ lái xe và xe tự hành: hệ thống phải xác định đúng loại biển trong ảnh từ camera để cảnh báo hoặc điều chỉnh hành vi của xe. Bài toán khó ở chỗ ảnh biển thay đổi nhiều về kích thước, độ sáng, góc nhìn, và nhiều biển khác nhau chỉ khác ở chi tiết nhỏ như chữ số trên biển giới hạn tốc độ.

Báo cáo này trình bày bài tập lớn học phần Nhập môn học máy, đề tài 6: phân loại ảnh biển báo giao thông vào 43 lớp trên bộ dữ liệu German Traffic Sign Recognition Benchmark (GTSRB) bằng các thuật toán đã học trong học phần là K-Nearest Neighbors (K-NN) và Cây quyết định. Nhóm xây dựng quy trình học máy hoàn chỉnh từ dữ liệu thô đến đánh giá, so sánh ba cách biểu diễn đặc trưng (pixel, HOG, PCA) và hai cách tiền xử lý (có và không có CLAHE), đồng thời phân tích các lớp dễ nhầm lẫn và các trường hợp dự đoán sai.

Báo cáo gồm ba chương. Chương 1 đặt vấn đề, phân tích yêu cầu và khám phá dữ liệu. Chương 2 trình bày thiết kế giải pháp, quy trình học máy và cấu trúc mã nguồn. Chương 3 trình bày thiết lập thí nghiệm, kết quả, phân tích lỗi, thảo luận và so sánh. Phần Kết luận tổng kết kết quả chính và hướng phát triển.
